\hypertarget{appendix-faac749c9cdd}{}
\section*{Half-edge topology invariants}
\chapter{Half-edge topology invariants}\label{appendix-faac749c9cdd-half-edge-topology-invariants}

A triangle index buffer tells the renderer which three vertices form each face, but many geometry-processing algorithms need explicit adjacency: which face lies across an edge, which edges surround a vertex, and where the mesh boundary is.

A half-edge representation splits every undirected edge into directed records attached to faces. For triangle face \((v_0,v_1,v_2)\), create directed half-edges

\[
v_0\rightarrow v_1,\quad
v_1\rightarrow v_2,\quad
v_2\rightarrow v_0.
\]

For a consistently oriented two-manifold interior edge, exactly two half-edges share the same undirected endpoints and they run in opposite directions. They become twins.

\section{Invariants}\label{appendix-faac749c9cdd-invariants}

A useful topology builder should establish:

\begin{itemize}
\tightlist
\item
  every face vertex index is valid;
\item
  a face does not repeat one of its own vertices;
\item
  each directed edge appears at most once for an oriented manifold triangle mesh;
\item
  an undirected edge with one incident face is a boundary edge;
\item
  an undirected edge with two incident faces is manifold only when directions oppose;
\item
  an undirected edge with more than two incident faces is non-manifold;
\item
  twin links are symmetric;
\item
  face-cycle next links return to the original half-edge.
\end{itemize}

Vertex-manifoldness needs an additional fan test. Merely checking that every edge has at most two incident faces does not reject a ``bow-tie'' vertex whose incident triangles form separate fans touching only at that vertex.

\section{Euler characteristic as a diagnostic}\label{appendix-faac749c9cdd-euler-characteristic-as-a-diagnostic}

For a finite mesh complex,

\[
\chi = V-E+F.
\]

A topological disk has \(\chi=1\). A closed sphere-like tetrahedron has \(\chi=2\). Euler characteristic does not fully classify arbitrary meshes, but it is a powerful cross-check against adjacency construction.

\section{Why topology must be separated from geometry}\label{appendix-faac749c9cdd-why-topology-must-be-separated-from-geometry}

Two faces can be topologically adjacent even if their coordinates are nearly coincident or wildly deformed. Conversely, two spatially overlapping triangles need not share topology. Build connectivity from explicit vertex/edge identity unless the task is specifically to weld geometric duplicates under a tolerance.

\section{Executable witness}\label{appendix-faac749c9cdd-executable-witness}

tools/labs/mesh\_topology\_lab.cpp builds directed/undirected edge maps, pairs twins, counts boundaries and non-manifold edges, checks orientation, and validates Euler characteristic on a two-triangle disk and a closed tetrahedron. A third case deliberately creates a three-face non-manifold edge.

The witness does not yet prove vertex-fan manifoldness, polygonal faces, robust welding, edge collapse, remeshing or production importer parity.

Related: \hyperref[appendix-1aea26b2307f]{Half-edge mesh connectivity}, \hyperref[appendix-e38976cbbf86]{Mesh topology and homology}, \hyperref[appendix-89db9fe1bc8a]{QEM edge collapse and validity}.
